% Part: first-order-logic
% Chapter: sequent-calculus
% Section: proof-theoretic-notions

\documentclass[../../../include/open-logic-section]{subfiles}

\begin{document}

\begin{editorial}
  يجمع هذا القسم تعريفي علاقة قابلية الإثبات والاتساق في حساب
  التتابعيات.
\end{editorial}

\iftag{FOL}
      {\olfileid{fol}{seq}{ptn}}
      {\olfileid{pl}{seq}{ptn}}

\olsection{مفاهيم نظرية البرهان}

\begin{explain}
مثلما عرّفنا عددًا من المفاهيم الدلالية المهمة (الصلاحية، والاستلزام،
وقابلية الإرضاء)، نعرّف الآن \emph{مفاهيم نظرية البرهان} المقابلة. ولا
تُعرّف هذه بالرجوع إلى إرضاء ال!!{sentence}s في
\iftag{FOL}{ال!!{structure}s}{ال!!{valuation}s}، بل
بالرجوع إلى !!{derivability} بعض التتابعيات أو
!!{nonderivability}ها. وكان اكتشاف تطابق هذه المفاهيم اكتشافًا مهمًا.
وتقرر هذا التطابق \emph{مبرهنتا السلامة} و\emph{الاكتمال}.
\end{explain}

\begin{defn}[المبرهنات]
تكون !!^a{sentence}~$!A$ \emph{مبرهنةً} إذا وجد !!a{derivation} في
~$\Log{LK}$ للتتابعية~$\quad \Sequent !A$. ونكتب $\Proves !A$ إذا
كانت~$!A$ مبرهنة، ونكتب~$\Proves/ !A$ إذا لم تكن كذلك.
\end{defn}

\begin{defn}[!!^{derivability}]
تكون !!^a{sentence}~$!A$ \emph{!!{derivable} من} مجموعة
!!{sentence}s~$\Gamma$، أي $\Gamma \Proves !A$، إذا وفقط إذا وجدت
مجموعة جزئية منتهية~$\Gamma_0 \subseteq \Gamma$ ومتتالية
$\Gamma_0'$ من ال!!{sentence}s في~$\Gamma_0$ بحيث
$\Log{LK}$ !!{derive}s التتابعية~$\Gamma_0' \Sequent !A$. وإذا لم
تكن~$!A$ !!{derivable} من~$\Gamma$، كتبنا $\Gamma \Proves/ !A$.
\end{defn}

بفضل قواعد الانكماش والإضعاف والتبادل، لا يهم ترتيب
ال!!{sentence}s في~$\Gamma_0'$ ولا عددها: فإذا كانت التتابعية
$\Gamma_0' \Sequent !A$ !!{derivable}، كانت
$\Gamma_0'' \Sequent !A$ كذلك لأي~$\Gamma_0''$ تحتوي على
ال!!{sentence}s نفسها التي تحتوي عليها~$\Gamma_0'$. فمثلًا، إذا كان
$\Gamma_0 = \{!B, !C\}$، كانت كل من
$\Gamma_0' = \tuple{!B, !B, !C}$ و
$\Gamma_0'' = \tuple{!C, !C, !B}$ متتالية لا تحتوي إلا على
ال!!{sentence}s في~$\Gamma_0$. وإذا كانت تتابعية تحتوي على إحداهما
!!{derivable}، كانت الأخرى كذلك، مثلًا:
\begin{prooftree}
  \AxiomC{}
  \Deduce$!B, !B, !C \fCenter !A$
  \RightLabel{\LeftR{\Contraction}}
  \UnaryInf$!B, !C \fCenter !A$
  \RightLabel{\LeftR{\Exchange}}
  \UnaryInf$!C, !B \fCenter !A$
  \RightLabel{\LeftR{\Weakening}}
  \UnaryInf$!C, !C, !B \fCenter !A$
\end{prooftree}
سنقول من الآن فصاعدًا إنه إذا كانت~$\Gamma_0$ مجموعة منتهية من
ال!!{sentence}s، فإن~$\Gamma_0 \Sequent !A$ هي أي تتابعية يكون
مقدَّمها متتاليةً من ال!!{sentence}s في~$\Gamma_0$، وسنضمّن ضمنًا
استدلالات الانكماش والتبادل والإضعاف عند الحاجة.

\begin{defn}[الاتساق]
تكون مجموعة الجمل~$\Gamma$ \emph{غير متسقة} إذا وفقط إذا وجدت مجموعة
جزئية منتهية~$\Gamma_0 \subseteq \Gamma$ بحيث
$\Log{LK}$ !!{derive}s التتابعية~$\Gamma_0 \Sequent \quad$. وإذا لم
تكن~$\Gamma$ غير متسقة، أي إذا كان، لكل~$\Gamma_0 \subseteq \Gamma$
منتهية، $\Log{LK}$ لا !!{derive} $\Gamma_0 \Sequent \quad$،
قلنا إنها \emph{متسقة}.
\end{defn}

\begin{prop}[الانعكاسية]
\ollabel{prop:reflexivity}
إذا كان~$!A \in \Gamma$، فإن~$\Gamma \Proves !A$.
\end{prop}

\begin{proof}
التتابعية الابتدائية~$!A \Sequent !A$ !!{derivable}، و
$\{!A\} \subseteq \Gamma$.
\end{proof}
  
\begin{prop}[الرتابة]
\ollabel{prop:monotonicity}
إذا كان~$\Gamma \subseteq \Delta$ و$\Gamma \Proves !A$، فإن
$\Delta \Proves !A$.
\end{prop}

\begin{proof}
افترض أن~$\Gamma \Proves !A$، أي إنه توجد مجموعة منتهية
$\Gamma_0 \subseteq \Gamma$ تكون~$\Gamma_0 \Sequent !A$
!!{derivable}. وبما أن~$\Gamma \subseteq \Delta$، فإن~$\Gamma_0$
مجموعة جزئية منتهية من~$\Delta$ أيضًا. ومن ثم يبين !!{derivation}
$\Gamma_0 \Sequent !A$ كذلك أن~$\Delta \Proves !A$.
\end{proof}

\begin{prop}[التعدي]
\ollabel{prop:transitivity}
إذا كان~$\Gamma \Proves !A$ و$\{!A\} \cup \Delta \Proves !B$، فإن
$\Gamma \cup \Delta \Proves !B$.
\end{prop}

\begin{proof}
إذا كان~$\Gamma \Proves !A$، وجدت مجموعة منتهية
$\Gamma_0 \subseteq \Gamma$ و!!a{derivation}~$\pi_0$ للتتابعية
$\Gamma_0 \Sequent !A$. وإذا كان~$\{!A\} \cup \Delta \Proves !B$،
فإنه توجد، لبعض مجموعة جزئية منتهية
$\Delta_0 \subseteq \Delta$، !!a{derivation}~$\pi_1$ للتتابعية
$!A, \Delta_0 \Sequent !B$. انظر في ال!!{derivation} الآتي:
\begin{prooftree}
  \AxiomC{}
  \RightLabel{$\pi_0$}
  \Deduce$\Gamma_0 \fCenter !A$
  \AxiomC{}
  \RightLabel{$\pi_1$}
  \Deduce$!A, \Delta_0 \fCenter !B$
  \RightLabel{\Cut}
  \BinaryInf$\Gamma_0, \Delta_0 \fCenter !B$
\end{prooftree}
بما أن~$\Gamma_0 \cup \Delta_0 \subseteq \Gamma \cup \Delta$، فإن
هذا يبين أن~$\Gamma \cup \Delta \Proves !B$.
\end{proof}

لاحظ أن هذا يعني خصوصًا أنه إذا كان~$\Gamma \Proves !A$ و
$!A \Proves !B$، فإن~$\Gamma \Proves !B$. ويلزم أيضًا أنه إذا كان
$!A_1, \dots, !A_n \Proves !B$ وكان~$\Gamma \Proves !A_i$ لكل~$i$،
فإن~$\Gamma \Proves !B$.

\begin{prop}
\ollabel{prop:incons}
تكون~$\Gamma$ غير متسقة إذا وفقط إذا كان~$\Gamma \Proves {!A}$ لكل
  جملة~$!A$.
\end{prop}

\begin{proof}
تمرين.
\end{proof}

\begin{prob}
أثبت \olref[fol][seq][ptn]{prop:incons}
\end{prob}

\begin{prop}[التراص]
  \ollabel{prop:proves-compact}
  \begin{enumerate}
  \item إذا كان~$\Gamma \Proves !A$، وجدت مجموعة جزئية منتهية
    $\Gamma_0 \subseteq \Gamma$ بحيث~$\Gamma_0 \Proves !A$.
  \item إذا كانت كل مجموعة جزئية منتهية من~$\Gamma$ متسقة، فإن
    $\Gamma$ متسقة.
  \end{enumerate}
\end{prop}

\begin{proof}
  \begin{enumerate}
    \item إذا كان~$\Gamma \Proves !A$، وجدت مجموعة جزئية منتهية
      $\Gamma_0 \subseteq \Gamma$ يكون للتتابعية
      $\Gamma_0 \Sequent !A$ !!a{derivation}. وبناءً على ذلك،
      $\Gamma_0 \Proves !A$.
    \item إذا كانت~$\Gamma$ غير متسقة، وجدت مجموعة جزئية منتهية
      ~$\Gamma_0 \subseteq \Gamma$ بحيث $\Log{LK}$ !!{derive}s
      التتابعية~$\Gamma_0 \Sequent \quad$. ولكن~$\Gamma_0$ عندئذ
      مجموعة جزئية منتهية من~$\Gamma$ وغير متسقة.
  \end{enumerate}
\end{proof}

\end{document}
